\documentclass[11pt]{article}

\usepackage[utf8]{inputenc}
\usepackage[T1]{fontenc}
\usepackage{lmodern}
\usepackage[a4paper,margin=1.05in]{geometry}
\usepackage{booktabs}
\usepackage{array}
\usepackage{parskip}
\usepackage[hidelinks]{hyperref}

\newcommand{\caveat}{\emph{\small Reminder: every median below is time-to-hit
among people who already succeeded. The sample has no denominator, so no figure
here supports a statement about the odds of succeeding.}\par\medskip}

\title{\bfseries Apprenticeship Trajectories\\[0.35em]
\large Interim status report at 160 people researched}
\author{\normalsize Generated from the study repository at \texttt{\textasciitilde/trajectories}}
\date{\normalsize 22 August 2026}

\begin{document}
\maketitle

\begin{abstract}
\noindent
This is a snapshot taken part way through data collection, not a result. As of
this writing the study has researched 160 people, included 134 of them, and can
compute its headline revenue-strict clock on 60 rows. On those 60 rows the
median time from the end of formal education to a first \$10M constant-2026
crossing is \textbf{9.2 years}, with a bootstrap 95\% confidence interval of
8.5 to 13.0 years. The study's own pre-registered stopping rule reports
\texttt{STOP: False}: neither the wave-over-wave stability criterion nor the
precision criterion is met. If collection continues as it has been going, the
figure the study is heading toward looks like roughly nine years of
apprenticeship. That is a description of where the running estimate sits, not a
finding. Above all, the median describes duration given success. It is
conditioned on the outcome, it has no denominator, and the two strongest
measured biases both push it short, so it should be read as a floor.
\end{abstract}

\section{Scope of this snapshot}

\textbf{These numbers cover waves through wave 6 only.} The pilot (10 people)
and waves 1 through 6 (25 people each) make 160 rows in
\texttt{data/anchors.csv}. \textbf{Waves 7 and 8, covering p161 through p210,
are drawn and in research right now, and they are not in any figure in this
document.} When those waves close, every number here moves. Nothing in this
report should be quoted without that qualification.

The state described is: 160 people researched, 134 included, 26 excluded, and
60 rows that satisfy the revenue-strict definition of the primary clock.

\section{What is being measured}

For each person the study records six dated career anchors: birth, end of
formal education, first full-time paid work in the field of the eventual hit,
first venture founded or independently authored, first hit, and the later scale
event they are famous for now. From these it computes four clocks. The headline
one is \texttt{clock\_education}, the number of years from
\texttt{a2\_education\_end} to \texttt{a5\_first\_hit}.

\subsection{The hit criterion}

For the commercial buckets, a hit is the first year the person's own venture
reached \textbf{\$10 million in constant 2026 US dollars}. The constant-dollar
part is load-bearing, and it was a correction made after the pilot.

A nominal \$10M threshold would have required a 1960 founder to build a
business roughly eleven times larger in real terms than a 2020 founder before
either counted as having hit. Since a large share of the sample has pre-1995
careers, a nominal bar would have inflated exactly those apprenticeships and
manufactured an era effect out of arithmetic. The deflated thresholds, from
\texttt{python -m src.cpi}, make the point concrete.

\begin{table}[h]
\centering
\begin{tabular}{lrrr}
\toprule
Revenue year & \texttt{rev10} & \texttt{acq50} & \texttt{fund100} \\
\midrule
1960 & \$919,417 & \$4,597,087 & \$9,194,174 \\
1995 & \$4,733,757 & \$23,668,786 & \$47,337,572 \\
2020 & \$8,039,032 & \$40,195,159 & \$80,390,318 \\
\bottomrule
\end{tabular}
\caption{Nominal dollar thresholds equivalent to the constant-2026 bars. The
1960 bar is about one eleventh of the 2026 bar.}
\end{table}

Non-USD accounts are converted at the spot rate for the revenue year, then
compared against that year's nominal threshold.

Where revenue is undisclosed, an IPO or an acquisition above \$50M constant is
permitted as a fallback basis, but only when no earlier crossing is known to
have occurred. A qualitative claim in a source, that the venture was already
large or industry-leading by a stated year, counts as establishing an earlier
crossing and bars the fallback. In that case the row gets a bounded date
\texttt{YYYY--YYYY} at most ten years wide, or it is excluded as
\texttt{crossing\_undatable}. This is deliberately strict: dating a row by a
later IPO produces a confidently sourced answer that is years wrong, which is
how the first pilot dated Gordon Moore about eleven years late.

The non-commercial buckets use per-bucket equivalents: a Nobel, Turing, Fields
or Breakthrough prize announcement for \texttt{science\_research}; a
million-person audience for \texttt{media\_creators}; a first \$100M constant
fund for fund managers, or a first year topping a recognised industry ranking
for analysts and economists; and revenue first, third-party volume rankings
second, for \texttt{trade\_import\_logistics}. These are not commensurable with
each other, which is precisely why the revenue-strict subset is the headline
and the pooled figure is reported only alongside it.

\section{Sample construction}

\subsection{Buckets}

The frame fixes eight field buckets with cumulative target shares. It was
frozen on 31 July 2026 and revised once, on 11 August 2026, after the
ten-person pilot. It has not been touched since wave 1 began, because changing
the frame mid-collection would invalidate the stopping rule: the sample
composition would no longer be stable across waves.

\begin{table}[h]
\centering
\begin{tabular}{lrrrrr}
\toprule
Bucket & Target & \multicolumn{2}{c}{Drawn (160)} & \multicolumn{2}{c}{Included (134)} \\
\cmidrule(lr){3-4}\cmidrule(lr){5-6}
 & share & $n$ & share & $n$ & share \\
\midrule
\texttt{software\_internet}           & 25\% & 40 & 25.0\% & 36 & 26.9\% \\
\texttt{hardware\_deeptech}           & 15\% & 24 & 15.0\% & 24 & 17.9\% \\
\texttt{consumer\_retail\_industrial} & 13\% & 21 & 13.1\% & 13 & 9.7\% \\
\texttt{science\_research}            & 12\% & 19 & 11.9\% & 19 & 14.2\% \\
\texttt{investors\_finance}           & 10\% & 16 & 10.0\% & 14 & 10.4\% \\
\texttt{media\_creators}              & 10\% & 16 & 10.0\% & 14 & 10.4\% \\
\texttt{healthcare\_biotech}          & 10\% & 16 & 10.0\% & 9  & 6.7\% \\
\texttt{trade\_import\_logistics}     & 5\%  & 8  & 5.0\%  & 5  & 3.7\% \\
\bottomrule
\end{tabular}
\caption{Bucket shares as drawn and as they survive exclusion. Drawing tracks
the target closely. Exclusion does not, and it pulls the analysed sample away
from the frame; see Section~\ref{sec:quality}.}
\end{table}

\subsection{Named source lists}

No name enters by free recall. Every roster entry cites the enumerated,
third-party list it came from, and roster membership must be evidenced by a URL
that names the person. Four names so far were credited to lists that turned out
not to contain them. The lists are: Forbes 400 (self-made score 6 to 10);
Forbes World's Billionaires (self-made); Midas List and Midas List Europe;
Fortune 40 Under 40; the Nobel Prizes in Physics, Chemistry, Medicine and
Economics; the Turing Award; the Fields Medal; the Breakthrough Prize; Time
100; Y Combinator Top Companies; the Hurun Rich List for China and India;
Nikkei and Toyo Keizai rankings for Japan; Maeil Business and Chosun Ilbo
rankings for Korea; the Sunday Times Rich List; the Pulitzer Prize; the Academy
Awards; the Grammy Awards; the Endeavor Entrepreneur network; and the Computer
History Museum Fellow Awards.

The last of these was added after wave 1 began. It was admitted on its own
merits, as an enumerated annual third-party honour list for computing with the
same shape as the Turing Award and Fields Medal already present, rather than to
admit any particular name.

\subsection{The selection rule}

Selection is money-only. The procedure that has worked, and that is
non-steering, is to enumerate a list in its published order and take the first
entries whose hit entity appears nowhere in \texttt{data/anchors.csv} or
\texttt{data/roster.csv}. Prominence is not a criterion in either direction.

That last clause was learned the hard way. A wave-6 roster agent was told to
prefer names it would not have recalled unprompted, and it duly rejected Peter
Thiel, Vinod Khosla, Jensen Huang and the Hinton/Hopfield pair from enumerated
lists as too obvious. Steering away from fame is a selection bias exactly as
much as recalling famous names is, and the two do not cancel: they widen the
variance and make the sample harder to characterise than either alone. This is
recorded as bias 24.

\subsection{Covariates, not quotas, and the wave-1 discontinuity}

Country, era and sex were originally enforced as floors: non-US at least 30\%,
pre-1995 hit at least 25\%, women at least 20\%, with any wave that would
breach a floor rebalanced by hand before research began. \textbf{After wave 1
the study author removed the floors.} The three variables are still recorded on
every row and reported every wave, but they no longer steer who is drawn.

The reasoning is that recording a variable lets you discover whether it moves
the answer, while selecting on it guarantees you cannot, because you fixed it
by hand. Era in particular is a genuine confounder for time-to-hit, since
capital availability, addressable market size and information flow all changed
around the internet era. Under quotas any drift was prevented and therefore
invisible. Under covariates it is measurable.

\textbf{This makes the sample a mixture of two sampling designs, and that must
be stated before any median is read.} The pilot and wave 1, rows 1 to 35, were
steered by hand and came in at 50\% pre-1995, 69\% non-US and 42\% women, well
above what the source lists produce unaided. Waves 2 onward were not steered.
The discontinuity is a step, not a gradient, and is recorded as bias 14. The
honest treatment once $N$ is large is to report the post-change rows separately
and check them against the pooled figure, not to retro-fit the early rows out:
they are correctly collected data, collected under a different rule.

The current cumulative composition of the 134 included rows is 86 non-US
(64.2\%), 45 women (33.6\%), and 26 pre-1995 hits (19.4\%) against 108
post-1995. If those shares drift hard, that is a finding to state, not a reason
to resume rebalancing.

\section{The current numbers}

\caveat

\subsection{Headline}

\begin{table}[h]
\centering
\begin{tabular}{lr}
\toprule
Revenue-strict $n$ & 60 \\
Median \texttt{clock\_education} & 9.2 yr \\
Bootstrap 95\% CI & 8.5 to 13.0 yr \\
CI half-width & 2.25 yr \\
Stopping-rule threshold on half-width & 1.00 yr \\
\bottomrule
\end{tabular}
\caption{Output of \texttt{python -m src.stats analysis/clocks.csv} at 160
people researched.}
\end{table}

Revenue-strict means the 60 included rows that carry hit basis
\texttt{primary}, that is an actual dated \$10M constant crossing rather than a
fallback or a per-bucket equivalent, and that also have a usable education
year. Of the 134 included rows, 77 carry basis \texttt{primary}, 49
\texttt{equivalent} and 8 \texttt{fallback}. The drop from 77 to 60 is missing
education years, discussed in Section~\ref{sec:quality}.

\subsection{Definition strictness}

\caveat

\begin{table}[h]
\centering
\begin{tabular}{lrl}
\toprule
Definition & $n$ & Median (95\% CI) \\
\midrule
Revenue-strict (\$10M revenue only)             & 60  & 9.2 yr (8.5 to 13.0) \\
All commercial (plus IPO/acquisition fallback)  & 68  & 10.0 yr (8.5 to 13.5) \\
Pooled (all buckets, mixed definitions)         & 103 & 14.0 yr (10.0 to 15.0) \\
\bottomrule
\end{tabular}
\caption{The headline and the pooled figure differ by 4.8 years. The pooled
number mixes prizes, audiences, funds and rankings with revenue crossings,
which is why it is not the headline.}
\end{table}

\subsection{The other three clocks}

\caveat

\begin{table}[h]
\centering
\begin{tabular}{lrl}
\toprule
Clock & $n$ & Median (95\% CI) \\
\midrule
\texttt{clock\_education}   & 103 & 14.0 yr (10.0 to 15.0) \\
\texttt{clock\_age18}       & 121 & 20.0 yr (18.0 to 23.0) \\
\texttt{clock\_venture}     & 114 & 5.0 yr (4.0 to 6.5) \\
\texttt{age\_at\_first\_hit}& 121 & 38.0 yr (36.0 to 41.0) \\
\bottomrule
\end{tabular}
\caption{All four clocks, pooled across every bucket and hit basis. These are
not the revenue-strict headline. They are printed together so a divergence
between the education clock and the age-18 clock stays visible, since that
divergence is the signal that the missing-education-year bias is biting.}
\end{table}

Note that \texttt{clock\_venture}, the years from first venture to first hit,
sits at 5.0 years while the education clock sits far above it. Most of the
measured apprenticeship is spent before the venture that hits exists.

\subsection{Wave-by-wave median history}

\caveat

\begin{table}[h]
\centering
\begin{tabular}{lrr}
\toprule
Wave closed & Revenue-strict $n$ & Median \\
\midrule
Pilot  & 13 & 10.0 yr \\
Wave 1 & 21 & 9.0 yr \\
Wave 2 & 28 & 8.8 yr \\
Wave 3 & 34 & 9.0 yr \\
Wave 4 & 45 & 9.0 yr \\
Wave 5 & 52 & 9.8 yr \\
Wave 6 & 60 & 9.2 yr \\
\bottomrule
\end{tabular}
\caption{Contents of \texttt{analysis/wave\_medians.txt}. The median has stayed
inside a 1.2-year band across seven closings while $n$ went from 13 to 60.}
\end{table}

This table is the basis for the claim in the abstract. The estimate has been
sitting near nine years for six consecutive closings, and its movement is not
growing. That is what "if the research keeps going as it has been, it looks
like roughly nine years" means, and it is the limit of what the history
supports. It describes where the running estimate has settled. It is not a
converged result, and the stopping rule below says so formally.

\section{The stopping rule}

The rule is pre-registered in \texttt{src/stats.py} and has three parts.

\begin{enumerate}
\item Revenue-strict $n$ of at least 30 (\texttt{N\_FLOOR}). \textbf{Met}, at
      $n = 60$. Below this floor the median is withheld entirely, in both the
      CLI and \texttt{analysis.md}.
\item Median drift under 0.50 years across two consecutive waves
      (\texttt{MAX\_DRIFT}). \textbf{Not met}: the last two drifts are 0.80 and
      0.60 years.
\item Bootstrap CI half-width at or under 1.00 years
      (\texttt{MAX\_HALF\_WIDTH}). \textbf{Not met}: it is 2.25 years.
\end{enumerate}

The CLI accordingly prints:

\begin{quote}
\texttt{STOP: False - median drift too large (0.80, 0.60 yr; need both < 0.50)}
\end{quote}

\subsection{Roughly what $n$ the rule implies}

Bootstrap median precision improves approximately with the square root of $n$.
Taking that at face value, closing the gap from a 2.25-year half-width to 1.00
year needs about
\[
n \approx 60 \times \left(\frac{2.25}{1.00}\right)^{2} \approx 304
\]
revenue-strict rows. At the current yield of 60 revenue-strict rows per 160
people researched, that is on the order of \textbf{800 people}, or roughly 26
further waves of 25.

Three qualifications on that arithmetic. First, the square-root rule is an
approximation and a bootstrap half-width for a median does not follow it
exactly, so the figure is an order of magnitude rather than a target. Second,
the yield ratio is not fixed: it depends on the exclusion rate and on the
missing education years discussed below, and improving either would lower the
people count without lowering the row count. Third, and most seriously, bias 23
means the half-width criterion can be satisfied on precision the study has not
earned; see Section~\ref{sec:limits}. Note that
\texttt{docs/HANDOFF.md} carries an older and lower estimate of this figure,
computed against a different half-width; the arithmetic above uses the current
2.25.

The rule refusing to call the number settled is correct behaviour, not a
failure. Stopping earlier and reporting the median with its actual interval,
plus the honest statement that it has not converged, is also a legitimate
scientific outcome. Which of those happens is a decision for the study author.

\section{Data quality}
\label{sec:quality}

\subsection{The exclusion audit}

Twenty-six of 160 rows are excluded, 16.2\%. Excluded rows appear in no other
table in the analysis, so a bias in what gets excluded is only visible here.

\begin{table}[h]
\centering
\begin{tabular}{lr}
\toprule
Exclusion reason & $n$ \\
\midrule
\texttt{crossing\_undatable}   & 22 \\
\texttt{roster\_unverified}    & 3 \\
\texttt{criterion\_never\_met} & 1 \\
\bottomrule
\end{tabular}
\caption{Why rows were dropped. An honest exclusion beats a fabricated date,
and every one of these has survived independent re-checking.}
\end{table}

\begin{table}[h]
\centering
\begin{tabular}{lrrr}
\toprule
Cross-cut & Included & Excluded & Exclusion rate \\
\midrule
US     & 48 & 8  & 14.3\% \\
non-US & 86 & 18 & 17.3\% \\
women  & 45 & 6  & 11.8\% \\
men    & 89 & 20 & 18.3\% \\
\midrule
\texttt{hardware\_deeptech}           & 24 & 0 & 0.0\% \\
\texttt{science\_research}            & 19 & 0 & 0.0\% \\
\texttt{software\_internet}           & 36 & 4 & 10.0\% \\
\texttt{investors\_finance}           & 14 & 2 & 12.5\% \\
\texttt{media\_creators}              & 14 & 2 & 12.5\% \\
\texttt{trade\_import\_logistics}     & 5  & 3 & 37.5\% \\
\texttt{consumer\_retail\_industrial} & 13 & 8 & 38.1\% \\
\texttt{healthcare\_biotech}          & 9  & 7 & 43.8\% \\
\bottomrule
\end{tabular}
\caption{Exclusion is not uniform. Two buckets lose nothing; two lose nearly
half. Era is not audited here, because an excluded row usually lacks the hit
year that era derives from, and judging a discarded row's era by hand would be
an inference rather than a measurement.}
\end{table}

The bucket pattern is the important one. \texttt{healthcare\_biotech} and
\texttt{consumer\_retail\_industrial} lose roughly two rows in five, almost
entirely to \texttt{crossing\_undatable}, because private companies in those
sectors do not publish early revenue. That is a whole-bucket failure mode
rather than scattered bad luck, and it is the same problem that once excluded
\texttt{trade\_import\_logistics} at 100\%, before the third-party volume
ranking criterion was added.

\subsection{Bounded dates}

\textbf{64 of the 134 included rows carry a bounded
\texttt{a5\_first\_hit}}, recorded as \texttt{YYYY--YYYY} with a source for
each end and a maximum span of ten years. Counting any bounded anchor rather
than \texttt{a5} alone, 65 of 134 rows are affected. The clocks use the
midpoint, and the report prints a sensitivity run so the uncertainty stays
visible.

\caveat

\begin{table}[h]
\centering
\begin{tabular}{lrl}
\toprule
Treatment & $n$ & Median (95\% CI) \\
\midrule
Bounded rows dropped        & 54  & 16.0 yr (13.0 to 19.0) \\
All rows, span midpoints    & 103 & 14.0 yr (10.0 to 15.0) \\
All rows, earliest possible & 103 & 12.0 yr (9.0 to 14.0) \\
All rows, latest possible   & 103 & 14.0 yr (12.0 to 18.0) \\
\bottomrule
\end{tabular}
\caption{Bounded-date sensitivity, pooled across all buckets and hit bases.
This run qualifies the pooled median, not the revenue-strict headline. The
envelope rows show the full range the data permits.}
\end{table}

A bounded date is a real measurement rather than a guess. \texttt{1960--1961}
says the sources place the event in those two years, whether from two sources
bracketing it or from one stating the range. A company's founding year is a
valid lower bound, since it cannot earn revenue before it exists. It is not
licence to widen a range until it contains a convenient year.

\subsection{The binding constraint is now \texttt{a2}, not \texttt{a5}}

\textbf{31 of the 134 included rows have
\texttt{a2\_education\_end = unknown}.} Those rows contribute nothing to the
headline clock, because the headline clock counts from the education year.
This, and not the datability of the hit, is now what limits the study: 77
included rows carry a \texttt{primary} hit basis, but only 60 reach the
revenue-strict median.

The loss is not random. When a person finished school is well documented for US
and UK figures and poorly documented for Indian and Chinese founders of the
1970s and 1980s, and for anyone whose path did not run through a credential
worth reporting. So the education clock is silently conditioned on people whose
schooling was recorded, which means Western, credentialed and recent, even
though selection into the study is money-only. A targeted second pass on an
earlier cohort of ten such rows resolved six, and all four that stayed unknown
were non-US.

\subsection{Confidence}

\caveat

96 of the 134 included rows fall below high confidence on at least one anchor.
Restricting the pooled clock to high-confidence rows only gives 21.0 years (95\%
CI 14.0 to 27.5) on $n = 38$, against 14.0 years on all 103. Those diverge by
7.0 years, which is the clearest signal in the dataset that the pooled median
cannot be read as a point estimate. Both rows of that comparison pool every
bucket and hit basis, so the run qualifies the pooled figure and not the
revenue-strict headline.

\section{Limitations}
\label{sec:limits}

\texttt{docs/BIASES.md} catalogues 24. The ones that bear on reading the
current number are below, each with the direction it pushes.

\begin{table}[h]
\centering
\begin{tabular}{p{0.05\textwidth}p{0.50\textwidth}p{0.30\textwidth}}
\toprule
\# & Bias & Direction \\
\midrule
1 & Survivorship. Everyone in the sample succeeded. Nobody appears who worked
    twelve or twenty years and never crossed. & Not a shift in value but a hard
    limit on meaning. Uncorrectable by more data. \\
\addlinespace
2 & Extreme-outcome selection, modest-milestone dating. Buckets are filled from
    lists that select the very top, but each person is dated at a milestone
    most passed decades earlier. People who ultimately reached Bezos scale
    probably crossed \$10M faster than people who crossed it and stopped there.
    & \textbf{Short.} \\
\addlinespace
3 & Selection into \texttt{primary} is non-random. A crossing is datable mainly
    when the company disclosed revenue early, which favours fast risers, and
    that is exactly the subset the stopping rule tracks. & \textbf{Short.} \\
\addlinespace
5 & Residual fallback lag on the 8 rows that legitimately use IPO or
    acquisition dating. & Long, on affected rows. \\
\addlinespace
6 & The \$10M level is a judgment, not a natural boundary. & Unknown. A lower
    bar would shorten every clock. \\
\addlinespace
7 & Non-commercial equivalents are not commensurable with revenue or with each
    other. & Unknown; adds variance to the pooled median only. \\
\addlinespace
14 & The sample mixes two sampling designs. Rows 1 to 35 were steered; rows 36
    onward were not. & Unsigned, and a step rather than a gradient. \\
\addlinespace
17 & The education clock is conditioned on documented schooling, and the 31
    missing years skew non-US and pre-1995. & Plausibly short. \\
\addlinespace
18 & Single-pass research has a measured miss rate. A rescue pass on 14
    excluded rows found 3 of 13 \texttt{crossing\_undatable} rows to be
    datable, and the recoveries skewed non-US and pre-1995. & Composition.
    Every standing exclusion is weaker evidence than it looks. \\
\addlinespace
22 & Blind-audit isolation is enforced by instruction, not by environment. Wave
    6's four concurrent research agents left 54 artifacts, including finished
    per-person answers, in a shared scratchpad. & If it ever fails it inflates
    apparent reliability, undetectably. \\
\addlinespace
23 & Co-founders of one company are not independent observations. p70 Dubinsky
    and p94 Hawkins both date to 1995 at Palm Computing: two people, one event.
    & \textbf{Falsely narrows the CI}, which attacks the stopping rule
    directly. \\
\addlinespace
24 & Steering away from fame is still steering. & Unknown, plausibly opposite
    to recall bias. Two opposing biases widen variance rather than cancelling. \\
\bottomrule
\end{tabular}
\caption{Selected biases and their directions, numbered as in
\texttt{docs/BIASES.md}.}
\end{table}

Two points deserve stating outside the table.

\textbf{The two strongest identified effects both push short.} Biases 2 and 3
work in the same direction, and nothing of comparable strength pushes the other
way. The working assumption should therefore be that the true median for a
normal successful career is longer than 9.2 years, and that the reported figure
is a floor.

\textbf{Bias 1 sits above all of it and cannot be corrected.} However precise
the figure becomes, it describes people who made it. The people who worked
twenty years and did not are absent by construction, not by rarity. A median of
9.2 years supports the sentence "among people who eventually hit, the middle
one took about 9.2 years." It does not support any sentence about the
probability of hitting, because that needs a denominator and this dataset has
none. Timing how long lottery winners held their tickets gives a real number,
correctly computed, that says nothing about whether buying tickets is wise.
$N = 1000$ would make the median precise and leave it exactly as uninformative
about odds.

Two mitigations that worked are worth recording alongside the failures. The
strongest reliability evidence the study has produced is a six-platform
cross-check, 60 comparisons with zero conflicts, on rows verified with no
access to the repository at all. And wave 6's blind audit came back at 0.00
disagreement on 10 drawn rows, with one second-pass-only miss.

\section{What happens next}

\begin{enumerate}
\item \textbf{Close waves 7 and 8}, p161 through p210, currently in research.
      That takes the sample to 210 people and will move every figure here.
\item \textbf{Re-run the pipeline} on the merged file: \texttt{src.clocks},
      \texttt{src.report}, then \texttt{src.stats --history}. Two more wave
      medians enter \texttt{analysis/wave\_medians.txt} and the drift criterion
      is re-evaluated against them.
\item \textbf{Attack \texttt{a2}, since it is the binding constraint.} A
      targeted second pass on the 31 included rows with an unknown education
      year would add rows to the headline clock without researching a single
      new person, and would specifically recover the non-US and pre-1995
      careers the clock is currently missing.
\item \textbf{Run a second independent pass on the 22
      \texttt{crossing\_undatable} rows.} The measured recovery rate on an
      earlier batch was roughly one in four, concentrated in non-English
      sources.
\item \textbf{Decide the Palm pair} (bias 23) and adopt a one-row-per-hit-event
      rule, entity plus year, rather than one row per person. That matters
      before the CI half-width is trusted, not after.
\item \textbf{Fix the blind-audit isolation} (bias 22) by running second passes
      in a clean working directory with no access to first-pass artifacts, and
      by carving the \texttt{data/anchors.csv} instruction out of any brief
      handed to an auditor.
\item \textbf{Resolve the seven items in \texttt{docs/OPEN-QUESTIONS.md}}, two
      of which decide whether a recorded row is included or excluded.
\end{enumerate}

Whether the study runs to the stopping rule's roughly 800 people, or stops
earlier and reports the interval as it stands, is the study author's call. Both
are defensible. Only one of them is currently the pre-registered default.

\vspace{1.5em}
\noindent\rule{\linewidth}{0.4pt}

\footnotesize
\noindent Every figure in this document was computed from the repository at 160
rows in \texttt{data/anchors.csv}. Medians and intervals come from
\texttt{python -m src.stats analysis/clocks.csv} and
\texttt{analysis/analysis.md}; wave history from
\texttt{analysis/wave\_medians.txt}; thresholds from \texttt{python -m src.cpi};
counts from \texttt{data/anchors.csv} and \texttt{analysis/clocks.csv}
directly. Bias numbering follows \texttt{docs/BIASES.md}.

\end{document}
